\documentclass{article}
\usepackage{graphicx}
\usepackage{amsmath}
\usepackage[T2A]{fontenc}  % For Cyrillic fonts
\usepackage[utf8]{inputenc}  % For UTF-8 encoding
\usepackage[russian]{babel}  % For Russian language support
\usepackage{listings}
\usepackage{xcolor}
\usepackage{mathtools}
\usepackage{color}

\definecolor{dkgreen}{rgb}{0,0.6,0}
\definecolor{gray}{rgb}{0.5,0.5,0.5}
\definecolor{mauve}{rgb}{0.58,0,0.82}

\lstset{frame=tb,
    language=Python,
    aboveskip=3mm,
    belowskip=3mm,
    showstringspaces=false,
    columns=flexible,
    basicstyle={\small\ttfamily},
    numbers=left,
    numberstyle=\tiny\color{gray},
    keywordstyle=\color{blue},
    commentstyle=\color{dkgreen},
    stringstyle=\color{mauve},
    breaklines=true,
    breakatwhitespace=true,
    tabsize=3
}
\title{Algorithms 1 \\ Homework 4}
\author{Nikita Zagainov}
\date{October 2024}  % Automatically display the current month and year

\begin{document}

\maketitle

\section{Sorting, Ordinal statistics}
\begin{enumerate}
    \item
          Чтобы эффективно находить количество уникальных
          элементов в объединении трех отсортированных
          массивов, можно хранить три указателя на каждый массив,
          а также по элементу, обозначающему последний элемент,
          встреченный в каждом массиве. На каждом шаге будем
          выбирать указатель, соответствующий минимальному
          из трех элементов, на который указывают все указатели,
          и проверять, соответствует ли этот элемент одному
          из последних элеиментов, которые мы храним.
          Если нет, то увеличиваем счётчик, далее в
          любом случае записываем этот элемент в последние
          элементы.

          \begin{lstlisting}[language=Python]
def count_unique_elems(arr1: list, arr2: list, arr3: list) -> int:
    arrs = [arr1, arr2, arr3]
    ptrs = [0, 0, 0]
    last_elems = [None, None, None]
    counter = 0

    while any(ptrs[i] < len(arrs[i]) for i in range(3)):
        to_incr = min(
            filter(lambda x: ptrs[x] < len(arrs[x]), range(3)),
            key=lambda x: arrs[x][ptrs[x]],
        )

        if arrs[to_incr][ptrs[to_incr]] not in last_elems:
            last_elems[to_incr] = arrs[to_incr][ptrs[to_incr]]
            counter += 1

        ptrs[to_incr] += 1

    return counter
    \end{lstlisting}

          Доказательство корректности:

          Каждый уникальный элемент будет добавлен хотя бы один раз,
          так как все три указателя полностью проходятся по
          массивам, и не добавить некий элемент алгоритм может
          только в случае, если такой элемент уже встречался.

          Ни один элемент не будет добавлен несколько раз.
          Рассмотрим какой-то шаг алгоритма:
          в $last\_elems$ указаны последние элементы из каждого
          массива. Мы выбираем указатель, соответствующий минимальному
          из трех элементов, на которые указывают все указатели.
          Если элемент в выбранном массиве уже встречался, то он
          гарантированно указан в одном из них. Противный случай ---
          это когда этот элемент уже был рассмотрен, но был перезаписан
          в $last\_elems$, но это противоречит алгоритму, так как он
          на каждом шаге выбирает указатель, соответствующий минимальному
          элементу.

          Оценка:

          Сложность этого алгоритма по времени --- $\Theta(n + m + k)$, так как
          мы один раз проходимся по каждому из массивов. Сложность алгоритма
          по памяти --- $\Theta(1)$, так как использование памяти не зависит от
          величины входных массивов.

    \item
          Заметим, что количество инверсий равно количеству раз,
          когда алгоритм сортировки меняет местами два элемента,
          стоящих рядом в неотсортированном порядке. Тогда единственная
          модификация в любой алгоритм сортировки для подсчёта инверсий ---
          это просто подсчёт таких попарных замен

          \begin{lstlisting}[language=Python]
def count_inversions(arr):
    temp_arr = arr.copy()
    return merge_sort_and_count(arr, temp_arr, 0, len(arr) - 1)


def merge_sort_and_count(arr, temp_arr, left, right):
    inv_count = 0
    if left < right:
        mid = (left + right) // 2

        inv_count += merge_sort_and_count(arr, temp_arr, left, mid)
        inv_count += merge_sort_and_count(arr, temp_arr, mid + 1, right)

        inv_count += merge_and_count(arr, temp_arr, left, mid, right)

    return inv_count


def merge_and_count(arr, temp_arr, left, mid, right):
    left_ptr = left
    right_ptr = mid + 1
    write_ptr = left
    inv_count = 0

    while left_ptr <= mid and right_ptr <= right:
        if arr[left_ptr] <= arr[right_ptr]:
            temp_arr[write_ptr] = arr[left_ptr]
            left_ptr += 1
        else:
            temp_arr[write_ptr] = arr[right_ptr]
            right_ptr += 1
            inv_count += mid - left_ptr + 1 # equivalent to pairwise swap of element on right_ptr with every element on its left
        write_ptr += 1

    while left_ptr <= mid:
        temp_arr[write_ptr] = arr[left_ptr]
        left_ptr += 1
        write_ptr += 1

    while right_ptr <= right:
        temp_arr[write_ptr] = arr[right_ptr]
        right_ptr += 1
        write_ptr += 1

    for left_ptr in range(left, right + 1):
        arr[left_ptr] = temp_arr[left_ptr]

    return inv_count
        \end{lstlisting}
          Доказательство:

          Этот алгоритм корректен по причинам, названным выше.
          Алгоритм представляет собой модификацию алгоритма сортировки
          с подсчётом парных замен.

          Оценка:

          Сложность этого алгоритма по времени:
          $\Theta(n \log n)$
          по мастер-теореме $T(n) = 2T(\frac{n}{2}) + n$.
          Сложность по памяти: $\Theta(n)$, так как данный алгоритм
          создаёт копию входного массива.

    \item
          Для решения этой задачи можно один раз проитерироваться
          по всем точкам, для каждой определяя количество прямых,
          покрывающих эту точку. Это можно сделать следующим образом:
          \begin{itemize}
              \item Отсортировать координаты концов каждого отрезка,
                    поставить им в соответствие флаг,
                    определяющий, правый или левый это конец
              \item Пробежаться по концам слева направо,
                    постоянно обновляя счётчик, показывающий
                    количество покрывающих отрезков
              \item Отрезки, подходящие под требования, записывать
                    в массив
          \end{itemize}

          \begin{lstlisting}[language=Python]
from math import ceil


def covered_point_set(segments: list[tuple[float, float]]) -> list[tuple[float, float]]:
    req_coverage = ceil(2 * len(segments) / 3)
    ends = sorted(
        [(segment[0], 1) for segment in segments]
        + [(segment[1], -1) for segment in segments]
    )
    answer = []

    cur_coverage = 0
    cur_segment_start = None

    for end in ends:
        cur_coverage += end[1]
        if cur_coverage == req_coverage:
            if cur_segment_start is None:
                cur_segment_start = end[0]

        elif cur_segment_start is not None:
            answer.append((cur_segment_start, end[0]))
            cur_segment_start = None

    if cur_segment_start is not None:
        answer.append((cur_segment_start, ends[-1][0]))

    return answer
          \end{lstlisting}

          Доказательство:

          Этот алгоритм корректен, так как для каждой точки
          множества он хранит точное число покрывающих ее отрезков.

          Оценка:

          Сложность этого алгоритма по времени: $\Theta(n \log n)$,
          так как он один раз проходится по каждому из $2n$ концов
          отрезков, предварительно их отсортировав. Сложность
          алгоритма по памяти: $\Theta(n)$, так как он создаёт
          дополнительный массив из концов, и его размер линейно
          зависит от размера входного массива.

    \item Решением данной задачи является медиана входного массива.
          Докажем, что это так: Если число $s = \mathrm{med}(a)$, тогда
          если уменьшить или увеличить $s$ на некую малую величину
          $\epsilon$, то целевая функция уменьшится на $\epsilon n$
          за счёт приближения $n$ элементам массива, и увеличится на
          $\epsilon (n + 1)$ за счёт удаления от $n + 1$ элементов.
          При увеличении $s$ на достаточно большое значение, чтобы
          оно соответствовало или превосходило элемент, больший
          медианы или уменьшении до элемента, меньшего медианы,
          целевая функция возрастёт больше, чем на $\epsilon$.

          Тогда задача решается поиском медианы с помощью черного
          ящика, и сложность этого алгоритма по времени: $\Theta(n)$.
\end{enumerate}
\end{document}
